% \appendix
\renewcommand{\thechapter}{C} % Set Appendix Chapter Label to D

\chapter{Time Measuring Methodology}
\label{app:time_measuring_methodology}

Due to Linux does not support reading CPU cycles directly using the rdcycle instruction, the \texttt{clock\_gettime} function (CLOCK\_MONOTONIC) is used to measure elapsed time in nanoseconds. The following C code snippet demonstrates how to read the CPU cycles using this method:

\begin{lstlisting}[language=C, caption={Time Measurement Function Used in Gemmini Tests}, label={lst:read_cycles_function}, captionpos=b]
// From generators/gemmini/software/gemmini-rocc-tests/include/gemmini_testutils.h
static uint64_t read_cycles() {
    struct timespec ts;
    if (clock_gettime(CLOCK_MONOTONIC, &ts) != 0) {
         perror("clock_gettime failed");
         return 0;
    }
    return (uint64_t)ts.tv_sec * 1000000000ULL + ts.tv_nsec;
}
\end{lstlisting}

Exclusive tests were performed to measure the overhead of this method, which was found to be negligible compared to the overall execution time of the tests. 

The result shows that the overhead of using \texttt{clock\_gettime(CLOCK\_MONO\allowbreak TONIC)} is the lowest, around 7.4 μs, while \texttt{gettimeofday} has a similar overhead of around 7.7 μs. The \texttt{clock\_gettime\allowbreak (CLOCK\_\allowbreak PRO\allowbreak CESS\_CPUTIME\_ID)} has a higher overhead of around 26.4 μs, which is expected as it provides more detailed CPU time information.

\begin{lstlisting}[language=C, caption={C Code for Timing Test}, label={lst:timing_test}, captionpos=b]
/*
 * timing_test.c - Test program to compare clock_gettime() vs gettimeofday()
 * 
 * This program tests the timing functions available on Linux and compares
 * their resolution, accuracy, and overhead.
 */

#define _POSIX_C_SOURCE 200809L  // Enable POSIX.1-2008 features
#define _DEFAULT_SOURCE          // Enable additional glibc features

#include <stdio.h>
#include <stdlib.h>
#include <time.h>
#include <sys/time.h>
#include <unistd.h>
#include <stdint.h>

// Test how many times we can call each timing function
#define NUM_ITERATIONS 10000

// Function to get time using clock_gettime(CLOCK_MONOTONIC)
uint64_t get_monotonic_time_ns() {
    struct timespec ts;
    if (clock_gettime(CLOCK_MONOTONIC, &ts) != 0) {
        perror("clock_gettime failed");
        return 0;
    }
    return (uint64_t)ts.tv_sec * 1000000000ULL + ts.tv_nsec;
}

// Function to get time using gettimeofday()
uint64_t get_timeofday_us() {
    struct timeval tv;
    if (gettimeofday(&tv, NULL) != 0) {
        perror("gettimeofday failed");
        return 0;
    }
    return (uint64_t)tv.tv_sec * 1000000ULL + tv.tv_usec;
}

// Function to get process CPU time
uint64_t get_process_time_ns() {
    struct timespec ts;
    if (clock_gettime(CLOCK_PROCESS_CPUTIME_ID, &ts) != 0) {
        perror("clock_gettime(CLOCK_PROCESS_CPUTIME_ID) failed");
        return 0;
    }
    return (uint64_t)ts.tv_sec * 1000000000ULL + ts.tv_nsec;
}

void test_timing_resolution() {
    printf("=== Testing Timing Resolution ===\n\n");
    
    // Test clock_gettime resolution
    struct timespec res;
    if (clock_getres(CLOCK_MONOTONIC, &res) == 0) {
        printf("CLOCK_MONOTONIC resolution: %ld.%09ld seconds\n", 
               res.tv_sec, res.tv_nsec);
        printf("  = %.1f nanoseconds\n", res.tv_sec * 1e9 + res.tv_nsec);
    } else {
        printf("Failed to get CLOCK_MONOTONIC resolution\n");
    }
    
    if (clock_getres(CLOCK_PROCESS_CPUTIME_ID, &res) == 0) {
        printf("CLOCK_PROCESS_CPUTIME_ID resolution: %ld.%09ld seconds\n", 
               res.tv_sec, res.tv_nsec);
        printf("  = %.1f nanoseconds\n", res.tv_sec * 1e9 + res.tv_nsec);
    } else {
        printf("Failed to get CLOCK_PROCESS_CPUTIME_ID resolution\n");
    }
    
    printf("gettimeofday() resolution: ~1 microsecond (1000 nanoseconds)\n");
    printf("\n");
}

void test_timing_accuracy() {
    printf("=== Testing Timing Accuracy (Minimal Operation) ===\n\n");
    
    uint64_t start, end, elapsed;
    volatile int dummy = 0;
    
    // Test clock_gettime(CLOCK_MONOTONIC)
    printf("Testing clock_gettime(CLOCK_MONOTONIC):\n");
    uint64_t min_mono = UINT64_MAX, max_mono = 0, total_mono = 0;
    
    for (int i = 0; i < 1000; i++) {
        start = get_monotonic_time_ns();
        dummy = dummy + 1;  // Minimal operation
        end = get_monotonic_time_ns();
        elapsed = end - start;
        
        if (elapsed < min_mono) min_mono = elapsed;
        if (elapsed > max_mono) max_mono = elapsed;
        total_mono += elapsed;
    }
    
    printf("  Min time: %lu ns\n", min_mono);
    printf("  Max time: %lu ns\n", max_mono);
    printf("  Avg time: %.1f ns\n", (double)total_mono / 1000.0);
    printf("  Effective resolution: ~%lu ns\n", min_mono);
    
    // Test gettimeofday()
    printf("\nTesting gettimeofday():\n");
    uint64_t min_gtod = UINT64_MAX, max_gtod = 0, total_gtod = 0;
    
    for (int i = 0; i < 1000; i++) {
        start = get_timeofday_us();
        dummy = dummy + 1;  // Minimal operation
        end = get_timeofday_us();
        elapsed = (end - start) * 1000;  // Convert to nanoseconds
        
        if (elapsed < min_gtod) min_gtod = elapsed;
        if (elapsed > max_gtod) max_gtod = elapsed;
        total_gtod += elapsed;
    }
    
    printf("  Min time: %lu ns\n", min_gtod);
    printf("  Max time: %lu ns\n", max_gtod);
    printf("  Avg time: %.1f ns\n", (double)total_gtod / 1000.0);
    printf("  Effective resolution: ~%lu ns\n", min_gtod);
    
    printf("\n");
}

void test_timing_overhead() {
    printf("=== Testing Timing Function Overhead ===\n\n");
    
    uint64_t start, end;
    uint64_t overhead_mono = 0, overhead_gtod = 0, overhead_proc = 0;
    
    // Test clock_gettime(CLOCK_MONOTONIC) overhead
    start = get_monotonic_time_ns();
    for (int i = 0; i < NUM_ITERATIONS; i++) {
        get_monotonic_time_ns();
    }
    end = get_monotonic_time_ns();
    overhead_mono = (end - start) / NUM_ITERATIONS;
    
    // Test gettimeofday() overhead  
    start = get_monotonic_time_ns();
    for (int i = 0; i < NUM_ITERATIONS; i++) {
        get_timeofday_us();
    }
    end = get_monotonic_time_ns();
    overhead_gtod = (end - start) / NUM_ITERATIONS;
    
    // Test clock_gettime(CLOCK_PROCESS_CPUTIME_ID) overhead
    start = get_monotonic_time_ns();
    for (int i = 0; i < NUM_ITERATIONS; i++) {
        get_process_time_ns();
    }
    end = get_monotonic_time_ns();
    overhead_proc = (end - start) / NUM_ITERATIONS;
    
    printf("Function call overhead (average over %d calls):\n", NUM_ITERATIONS);
    printf("  clock_gettime(CLOCK_MONOTONIC):     %lu ns\n", overhead_mono);
    printf("  gettimeofday():                     %lu ns\n", overhead_gtod);
    printf("  clock_gettime(CLOCK_PROCESS_CPUTIME): %lu ns\n", overhead_proc);
    printf("\n");
}

void test_timing_demo() {
    printf("=== Practical Timing Example ===\n\n");
    
    uint64_t start_mono, end_mono;
    uint64_t start_gtod, end_gtod;
    uint64_t start_proc, end_proc;
    
    printf("Timing a simple computation (1000 additions):\n");
    
    // Time with clock_gettime(CLOCK_MONOTONIC)
    start_mono = get_monotonic_time_ns();
    volatile long sum = 0;
    for (int i = 0; i < 1000; i++) {
        sum += i;
    }
    end_mono = get_monotonic_time_ns();
    
    // Time with gettimeofday()
    start_gtod = get_timeofday_us();
    sum = 0;
    for (int i = 0; i < 1000; i++) {
        sum += i;
    }
    end_gtod = get_timeofday_us();
    
    // Time with process CPU time
    start_proc = get_process_time_ns();
    sum = 0;
    for (int i = 0; i < 1000; i++) {
        sum += i;
    }
    end_proc = get_process_time_ns();
    
    printf("Results:\n");
    printf("  clock_gettime(CLOCK_MONOTONIC):     %lu ns (%.2f μs)\n", 
           end_mono - start_mono, (end_mono - start_mono) / 1000.0);
    printf("  gettimeofday():                     %lu μs (%.2f μs)\n", 
           end_gtod - start_gtod, (double)(end_gtod - start_gtod));
    printf("  clock_gettime(CLOCK_PROCESS_CPUTIME): %lu ns (%.2f μs)\n", 
           end_proc - start_proc, (end_proc - start_proc) / 1000.0);
    printf("  Sum result: %ld\n", sum);
    printf("\n");
}

void show_current_time() {
    printf("=== Current Time Values ===\n\n");
    
    struct timespec ts;
    struct timeval tv;
    
    // Show CLOCK_MONOTONIC
    if (clock_gettime(CLOCK_MONOTONIC, &ts) == 0) {
        printf("CLOCK_MONOTONIC: %ld.%09ld seconds\n", ts.tv_sec, ts.tv_nsec);
    }
    
    // Show gettimeofday
    if (gettimeofday(&tv, NULL) == 0) {
        printf("gettimeofday():  %ld.%06ld seconds\n", tv.tv_sec, tv.tv_usec);
    }
    
    // Show CLOCK_PROCESS_CPUTIME_ID
    if (clock_gettime(CLOCK_PROCESS_CPUTIME_ID, &ts) == 0) {
        printf("PROCESS_CPUTIME: %ld.%09ld seconds\n", ts.tv_sec, ts.tv_nsec);
    }
    
    printf("\n");
}

int main() {
    printf("Linux Timing Functions Test Program\n");
    printf("====================================\n\n");
    
    printf("This program tests the timing functions used in the Gemmini\n");
    printf("Linux port to replace the rdcycle instruction.\n\n");
    
    show_current_time();
    test_timing_resolution();
    test_timing_accuracy();
    test_timing_overhead();
    test_timing_demo();
    
    return 0;
}
\end{lstlisting}

\begin{lstlisting}[language=bash, caption={Output of Timing Test}, label={lst:timing_test_output}, captionpos=b]
debian@debian:~$ ./timing_test
Linux Timing Functions Test Program
====================================

This program tests the timing functions used in the Gemmini
Linux port to replace the rdcycle instruction.

=== Current Time Values ===

CLOCK_MONOTONIC: 288.119331000 seconds
gettimeofday():  1710323247.363538 seconds
PROCESS_CPUTIME: 0.039620000 seconds

=== Testing Timing Resolution ===

CLOCK_MONOTONIC resolution: 0.000000001 seconds
  = 1.0 nanoseconds
CLOCK_PROCESS_CPUTIME_ID resolution: 0.000000001 seconds
  = 1.0 nanoseconds
gettimeofday() resolution: ~1 microsecond (1000 nanoseconds)

=== Testing Timing Accuracy (Minimal Operation) ===

Testing clock_gettime(CLOCK_MONOTONIC):
  Min time: 6000 ns
  Max time: 133000 ns
  Avg time: 7505.0 ns
  Effective resolution: ~6000 ns

Testing gettimeofday():
  Min time: 7000 ns
  Max time: 133000 ns
  Avg time: 8165.0 ns
  Effective resolution: ~7000 ns

=== Testing Timing Function Overhead ===

Function call overhead (average over 10000 calls):
  clock_gettime(CLOCK_MONOTONIC):     7393 ns
  gettimeofday():                     7666 ns
  clock_gettime(CLOCK_PROCESS_CPUTIME): 26380 ns

=== Practical Timing Example ===

Timing a simple computation (1000 additions):
Results:
  clock_gettime(CLOCK_MONOTONIC):     69000 ns (69.00 μs)
  gettimeofday():                     69 μs (69.00 μs)
  clock_gettime(CLOCK_PROCESS_CPUTIME): 128000 ns (128.00 μs)
  Sum result: 499500
\end{lstlisting}
